\section{Conclusion}\label{sec-conclusion}

Nous avons montré qu'un fine-tuning modeste (3\,000 images, modèle
nano, 50 epochs sur un GPU T4) suffit à transformer un détecteur
générique en un détecteur adapté au trafic urbain dense, et que le gain
se concentre là où il était attendu : le rappel des véhicules de moins
de $32 \times 32$ pixels passe de $0{,}122$ à $0{,}694$
($+469{,}4\,\%$), tandis que la mAP@0,5 progresse de $0{,}438$ à
$0{,}825$ ($+88{,}4\,\%$), à débit d'inférence au moins égal. La valeur
du travail tient autant au protocole qu'aux chiffres : classes alignées
sur COCO et réindexées pour une comparaison équitable, rappel ventilé
par taille avec aires normalisées à la résolution d'entrée, jeu de test
isolé et dérivé de façon déterministe.

Trois prolongements se dégagent, par ordre de priorité. D'abord, la
constitution d'un jeu de données annoté de trafic béninois, préalable
indispensable pour transformer un résultat obtenu sur un domaine
comparable en un résultat établi sur le domaine cible, et dont ce
travail fournit le protocole d'évaluation clef en main. Ensuite, le
passage à l'échelle : entraîner sur les 45\,000 images de BMD-45 et sur
des déclinaisons plus grandes de YOLOv8 (s, m), pour mesurer ce que le
plafond observé sur les petits objets doit au modèle et ce qu'il doit à
la limite informationnelle. Enfin, la poursuite du pipeline de
vidéo-verbalisation (suivi multi-objets puis lecture de plaques), dont
la détection fiabilisée par ce travail constitue le socle.
